\section{Implementation And Evaluation}\label{sec:impl}
%
We have implemented a tool based on the above approach, called
\name{}, that synthesizes effectful test generators from user-provided
\tyName{} specifications and property.  The output program in our DSL
can be transpiled to OCaml (QCheck) or P.  Our experiments use the
OCaml transpilation to test effectful abstract datatype libraries,
serializers, transformers, and interpreters, concurrent data
structures and databases; we use the P target to test reactive
distributed system models (i.e., message-passing systems defined as a
collection of actors) by using the synthesized generator as a
property-directed scheduler that replaces P's default random one.
Notably, the same specification and synthesis strategy is used in both
cases.  \name{} consists of approximately 14K lines of OCaml code and
uses Z3~\cite{de2008z3} as its backend SMT solver.\footnote{The
  supplemental material provides detail results of our
  experiments. % as well as a docker image that contains the \name{}
  % source code and benchmarks.
}

\paragraph{Evaluation setting} \name{} takes two inputs: a target
property, expressed as an SRE, and an operator context that contains
\tyName{} specifications for the operators used in the SUT.  During
synthesis, \name{} decomposes the target property to build a generator
consistent with the traces that the property accepts, but constrained
to respect the history and future traces of the any effects it
references.  Because operator types are interpreted
underapproximately, there are potentially many valid candidate
programs that are consistent with the \tyName{} specifications of the
operators they use.  Our implementation considers a union of
well-typed synthesis candidates, the number of which is bounded by the
amount of time allotted for synthesis (3 minutes in our experiments).
% Non-determinism also manifests in how recursive functions are unrolled
% - recall that in an underapproximate setting, the body of a recursive
% function need only be unrolled until it can be shown that the effects
% performed by the loop's continuation are justified by the effects
% performed within the loop body and its history trace.
Synthesized generators use $\mykw{assume}$ statements to ensure that
data dependencies among effectful operations are respected and use
$\mykw{assert}$ statements to flag property violations.

%% We evaluate our approach by integrating our synthesized controllers
%% into the testing framework provided by P~\cite{DGJ+13,DQ17}, a
%% state-machine based, actor-style programming language tailored for
%% modeling distributed systems and testing user-defined safety and
%% liveness properties.
%% % implementing the operational semantics in
%% % \autoref{fig:selected-semantics}, where the auxiliary handler
%% % semantics ($\alpha \vDash \zevent{op}{\overline{c}} \Downarrow \beta$)
%% % can be implemented as a random sampler to provide messages $\beta$
%% % that align with handler specifications. However, to evaluate our
%% % approach in a more realistic setting,
%% In the P framework, actors are executable programs that communicate
%% via message passing.  To test a system, P's runtime monitors message
%% traffic between actors, checking that global safety and liveness
%% properties are maintained. By default, P’s runtime scheduler
%% systematically explores arbitrary message interleavings during
%% execution.

%% To test our synthesized controller with P handlers while also
%% retaining scheduling control, we translate our controllers into a
%% special P component that coordinates the messages between the actors
%% under test. In this setup, all messages are routed to our controller,
%% where they are buffered and then forwarded to the actual actors
%% according to the order found in the synthesized output.  The order in
%% which messages are forwarded from the controller is determined by
%% $\mykw{obs}$ statements in the controller program, allowing it to
%% control message scheduling.  The coordinator is also responsible for
%% generating and sending messages from the environment (e.g., a logical
%% user) to the actors under test, again respecting the order in which
%% these messages appear in the synthesized program.  The assume
%% statements in the translated coordinator ensure that generated
%% messages always have the expected payloads; assertion failures
%% indicate that the system under test did not encounter the potential
%% bug in this execution, indicating the need for another attempt.
%% % \BD{Say something about how we handle assertion failures.}
Our experimental evaluation addresses three research questions:
\begin{itemize}
\item[\textbf{Q1}:] Is \name{} \emph{expressive}? Can it synthesize
  generators for a diverse set of benchmarks with interesting
  non-trivial structural and semantic safety properties?
\item[\textbf{Q2}:] Is \name{} \emph{effective}? Do synthesized
  generators enable more targeted exploration of the state
  space to witness violations of provided safety properties than
  existing techniques?
\item[\textbf{Q3}:] Is \name{} \emph{efficient}? Is it able to
  synthesize meaningful generators in a reasonable amount
  of time?
\end{itemize} %

\newcolumntype{P}[1]{>{\centering\arraybackslash}p{#1}}

\begin{table}[t!]
  \renewcommand{\arraystretch}{0.8}
  % \vspace*{-.1in}
  \caption{\small Using \name{} to synthesize effectful test
    generators in OCaml. Benchmarks from prior work are annotated with
    their source: HAT~\cite{ZYDJ24}, the OCaml multicore PBT
    framework~\cite{OcamlMulticorePBT}, QuickCheck on IFC
    machines~\cite{pbt-ifc} and well-typed De Bruijn STLC
    programs~\cite{CoverageType}; as well as
    OLTP-Bench~\cite{OLTPBench}, an extensible testbed for
    benchmarking relational databases, and MonkeyDB~\cite{MonkeyDB}, a
    mock storage system for testing weak isolation levels.  The
    baselines used for comparison are taken from the original sources
    when they exist (these are labeled with $^{\dagger}$), otherwise
    we use QCheck's default random generator.  \name{} synthesized
    correct-by-construction test generators for all benchmarks within
    3 minutes; the number of generator variants synthesized is bounded
    at 3.  We set a 30-minute bound for the baselines to find a
    property violation.}
  \vspace*{-.1in} \footnotesize
%\multicolumn{2}{c||}{} &
\setlength{\tabcolsep}{1.3pt}
\begin{tabular}{c|P{4.8cm}||ccc||cc||c|ccc}
  \toprule
  \multicolumn{2}{c||}{Benchmark} & \#$\eff{op}$ & \multicolumn{2}{c||}{\#qualifier} & \multicolumn{2}{c||}{\scriptsize \# Num. Executions} & t$_\text{total}$(s) & \#evt & \#refine & \#SMT \\
  \midrule
Name & Property &   &uHAT &goal & {\scriptsize Clouseau} & {\scriptsize Baseline} &  &  \\
\midrule
\textsf{Stack}\cite{ZYDJ24} & {\scriptsize  Pushes and pops are correctly paired.} & $6$ & $12$ & $3$ & $1.0$ & $31.9$ & $0.60$ & $15$ & $4$ & $98$\\
\midrule
\textsf{Set}\cite{ZYDJ24} & {\scriptsize Membership holds for every element inserted into the set.} & $7$ & $22$ & $3$ & $1.0$ & $22.1$ & $1.28$ & $12$ & $12$ & $255$\\
\midrule
\textsf{Filesystem}\cite{ZYDJ24} & {\scriptsize A valid file path only contains non-deleted entries.} & $7$ & $67$ & $4$ & $1.3$ & $2812.1$ & $6.02$ & $13$ & $14$ & $853$\\
\midrule
\midrule
\textsf{Graph}\cite{ZYDJ24} & {\scriptsize A serialized stream of nodes and edges is re-constituted to
  form a fully-connected graph.} & $6$ & $24$ & $4$ & $1.0$ & {\tiny Timeout} & $16.61$ & $16$ & $16$ & $2114$\\
\midrule
\textsf{NFA}\cite{ZYDJ24} & {\scriptsize An NFA reaches a final state for every string in the language it accepts.} & $8$ & $50$ & $4$ & $1.0$ & {\tiny Timeout} & $21.85$ & $26$ & $6$ & $5080$\\
\midrule
\textsf{IFCStore}\cite{pbt-ifc} & {\scriptsize A well-behaved IFC program containing a $\Code{Store}$ command never leaks a secret.} & $8$ & $37$ & $2$ & $2.8$ & $4811.7^{\dagger}$ & $1.02$ & $8$ & $8$ & $160$\\
\midrule
\textsf{IFCAdd}\cite{pbt-ifc} & {\scriptsize A well-behaved IFC program containing an $\Code{Add}$ command never leaks a secret.} & $8$ & $37$ & $2$ & $1.6$ & $3383.4^{\dagger}$ & $0.96$ & $8$ & $8$ & $154$\\
\midrule
\textsf{IFCLoad}\cite{pbt-ifc} & {\scriptsize A well-behaved IFC program containing a $\Code{Load}$ command never leaks a secret.} & $8$ & $37$ & $2$ & $15.4$ & $11293.7^{\dagger}$ & $5.76$ & $16$ & $18$ & $930$\\
\midrule
\textsf{DeBruijn1}\cite{CoverageType} & {\scriptsize An STLC interpreter correctly evaluates a
  well-typed first-order STLC program that uses a de Brujin representation.} & $10$ & $117$ & $4$ & $1.5$ & $634.0^{\dagger}$ & $40.38$ & $18$ & $26$ & $4765$\\
\midrule
\textsf{DeBruijn2}\cite{CoverageType} & {\scriptsize An STLC interpreter correctly evaluates a
  a well-typed higher-order STLC program that uses a de Brujin representation.} & $10$ & $118$ & $4$ & $1.4$ & {\tiny Timeout}$^{\dagger}$ & $91.73$ & $18$ & $31$ & $9034$\\
\midrule
\midrule
\textsf{Shopping}\cite{MonkeyDB} & {\scriptsize All items added to a cart can be checked-out.} & $10$ & $60$ & $3$ & $1.0$ & $20.0^{\dagger}$ & $22.10$ & $17$ & $30$ & $1269$\\
\midrule
\textsf{HashTable}\cite{OcamlMulticorePBT} & {\scriptsize No updates to a concurrent hashtable are every lost.} & $13$ & $38$ & $4$ & $1.0$ & $2.5^{\dagger}$ & $0.76$ & $6$ & $6$ & $57$\\
\midrule
\textsf{Transaction} & {\scriptsize Asynchronous read operations are logically atomic.} & $8$ & $45$ & $3$ & $1.2$ & {\tiny Timeout} & $2.89$ & $15$ & $16$ & $423$\\
\midrule
\textsf{Courseware}\cite{MonkeyDB} & {\scriptsize Every student enrolled in a course exists in the
  enrollment database for that course.} & $16$ & $106$ & $3$ & $1.0$ & $57.5^{\dagger}$ & $27.71$ & $18$ & $33$ & $1479$\\
\midrule
\textsf{Twitter}\cite{MonkeyDB} & {\scriptsize Posted tweets are visible to all followers.} & $16$ & $99$ & $3$ & $1.0$ & $6.3^{\dagger}$ & $61.96$ & $24$ & $24$ & $2339$\\
\midrule
\textsf{Smallbank}\cite{OLTPBench} & {\scriptsize Account updates are strongly consistent.} & $22$ & $162$ & $3$ & $1.9$ & {\tiny Timeout} & $163.55$ & $28$ & $29$ & $10263$\\
\bottomrule
\end{tabular}
\label{tab:evaluation}
\vspace*{-.15in}
\end{table}



\begin{table}[t!]
  \renewcommand{\arraystretch}{0.8}
  % \vspace*{-.1in}
  \caption{\small Experimental results of using \name{} to synthesize
    property-directed schedulers for reactive distributed
    systems. Benchmarks from prior work are annotated with their
    source: P~\cite{DGJ+13}($^{\dagger}$),
    ModP~\cite{ModP}($^{\diamond}$) an extension of P with support for
    modules, and MessageChain~\cite{MessageChain}($^{\star}$), an
    automated verification tool for P.  We also include a real-world
    model of a two-phase commit protocol
    (\textsf{AnonReadAtomicity}$^{\square}$) used by a major cloud vendor.
    The components under test are written in P, and handler
    specifications for the SUTs components (message-passing actors)
    are given as \tyName{}s.  \name{} synthesizes a set
    of schedulers, each of which specifies a distinct scheduling
    order for messages, all consistent with provided specifications.
    We set a $3$ minute time bound for the synthesis procedure
    (t$_\text{total}$ is the average time to find a single
    controller.)  We set a bound of 30 minute time bound for the P
    baselines to generate a faulty execution.}
\vspace*{-.1in}
\footnotesize
\setlength{\tabcolsep}{1.3pt}
\begin{tabular}{c|P{4.0cm}||ccc||cc||c|ccc}
  \toprule
  \multicolumn{2}{c||}{Benchmark} & \#$\eff{op}$ & \multicolumn{2}{c||}{\#qualifier} & \multicolumn{2}{c||}{\scriptsize \# Num. Executions} & t$_\text{total}$(s) & \#evt & \#refine & \#SMT \\
  \midrule
Name & Property &   &uHAT &goal & {\scriptsize Clouseau} & {\scriptsize Baseline} &  &  \\
  \midrule
  {\scriptsize\textsf{Database}} & {\scriptsize The database maintains a Read-Your-Writes policy.} & $4$ & $15$ & $3$ & $1.0$ & $5.6$ & $0.32$ & $6$ & $6$ & $49$\\
  \midrule
  {\scriptsize\textsf{Firewall}}\cite{MessageChain} & {\scriptsize  Requests generated inside the firewall are eventually propagated to the outside.} & $5$ & $26$ & $4$ & $1.0$ & $10.0$ & $0.59$ & $10$ & $10$ & $222$\\
  \midrule
  {\scriptsize\textsf{RingLeaderElection}}\cite{MessageChain} & {\scriptsize There is always a single unique  leader.} & $3$ & $14$ & $2$ & $1.0$ & $17.7$ & $0.83$ & $8$ & $8$ & $100$\\
  \midrule
  {\scriptsize\textsf{BankServer}}\cite{DGJ+13} & {\scriptsize Withdrawals in excess of the available balance are never allowed.} & $6$ & $42$ & $1$ & $1.0$ & $2.2^{\dagger}$ & $0.17$ & $5$ & $5$ & $27$\\
  \midrule
  {\scriptsize\textsf{Simplified2PC}}\cite{DGJ+13} & {\scriptsize Transactions are atomic.} & $9$ & $32$ & $5$ & $2.1$ & $5.8^{\dagger}$ & $1.03$ & $8$ & $8$ & $88$\\
  \midrule
  {\scriptsize\textsf{HeartBeat}}\cite{DGJ+13} & {\scriptsize All available nodes are identified by a  detector.} & $7$ & $31$ & $3$ & $1.0$ & $6.0^{\dagger}$ & $1.10$ & $14$ & $14$ & $145$\\
  \midrule
  {\scriptsize\textsf{ChainReplication}}\cite{ModP} & {\scriptsize Concurrent updates are never lost.} & $7$ & $72$ & $4$ & $1.0$ & $500.0^{\dagger}$ & $19.24$ & $12$ & $169$ & $2054$\\
  \midrule
  {\scriptsize\textsf{Paxos}}\cite{ModP} & {\scriptsize Logs are correctly replicated.} & $10$ & $110$ & $2$ & $1.0$ & $667.7^{\dagger}$ & $23.70$ & $14$ & $77$ & $1763$\\
  \midrule
  {\scriptsize\textsf{Raft}} & {\scriptsize Leader election is robust to faults.} & $9$ & $39$ & $6$ & $1.0$ & {\tiny Timeout} & $26.84$ & $12$ & $78$ & $1262$\\
  \midrule
  {\scriptsize\textsf{AnonReadAtomicity}} & {\scriptsize Read Atomicity is preserved.} & $17$ & $113$ & $3$ & $1.0$ & $53.3^{\dagger}$ & $18.35$ & $16$ & $16$ & $1909$\\
\bottomrule
  \end{tabular}
\label{tab:evaluation2}
\vspace*{-.15in}
\end{table}

We have evaluated \name{} on a diverse corpus of programs drawn from a
variety of sources (described in the captions of
\autoref{tab:evaluation} and~\autoref{tab:evaluation2}); all of the
benchmarks except for \textsf{Database} and \textsf{Raft} were written
by others (\textbf{Q1}).  Our benchmarks are grouped into four
categories.  The first are data structure benchmarks (\textsf{Stack},
\textsf{Set}, \textsf{Filesystem}) where the test generator probes for
violations of structural invariants.  The second are evaluators -
programs that input serialized data, transform them into an internal
AST representation, and then interpret them; we consider evaluators
for data structures like graphs and automata (\textsf{Graph},
\textsf{NFA}), information flow control machines, and simply-typed
lambda-calculus, the latter two both described
in~\autoref{sec:overview}.  The third are concurrent and database
benchmarks (\textsf{Transaction}, \textsf{Shopping},
\textsf{Courseware}, \textsf{Twitter}, \textsf{Smallbank}) that test
properties related to consistency and synchronization. 
Property violations in any of these categories only manifest in specific configurations
(e.g., a violating test case for \textsf{NFA} involves constructing an
automata containing multiple nodes with more than two outgoing edges).  We
evaluate these three classes of benchmarks by transpiling synthesized
generators into OCaml, leveraging its QCheck library as
needed. The SUT for each of these benchmarks is known to be buggy; in
all cases we use implementations either available from prior sources
directly (if the OCaml source was available), or by directly
translating them into OCaml.  The fourth group of benchmarks shown
in~\autoref{tab:evaluation2} considers applications involving various
kinds of distributed protocols; these are executed using the P runtime
environment, by transpiling synthesized generators written in our DSL
into P's state machine modeling language.  Here as well, the protocol
being tested is known to be faulty from the sources they were adopted
from.


%% For example, \textsc{RingLeaderElection}
%% models a leader election algorithm.  Each message transmitted by the
%% model's actors contain a source ($\I{src}$) and destination
%% ($\I{dest}$), to allow recipient's to distinguish
%% messages from different nodes. Then, an $\eff{elect}$ handler can
%% be given the following \PAT{}: {\small\begin{align*}
%%   &\Code{n}\hasoty{tNode}{\top} \sarr \Code{m}\hasoty{tNode}{\top}
%%   \sarr \Code{ld}\hasoty{tNode}{v = \Code{m}} \sarr
%%   \\&\rg{\globalA\evparenth{\top}}{\lastA\msgB{elect}{\I{src}\;\I{dest}\;\I{leader}}{
%%       \I{src} = \Code{n} \land \I{dest} = \Code{m} \land \I{leader} =
%%       \Code{m} }}{\finalA\msgB{beLeader}{\I{leader}}{\I{leader} =
%%       \Code{m}}}
%% \end{align*}}\noindent
%% This specification captures a useful property of a leader
%% election algorithm: when a node $\Code{m}$ receive a message that elects
%% ($\eff{elect}$) itself as leader ($\I{leader} = \Code{m}$) from
%% another node $\Code{n}$, then node $\Code{m}$ can announce itself to
%% be the leader ($\eff{beLeader}$).

\autoref{tab:evaluation} divides the results of our experiments into
four categories, separated by double bars. The first measures the
complexity of benchmarks with respect to the number of distinct
effectful operators (\#$\eff{op}$) and the number of qualifiers (i.e.,
symbolic events) used in \tyName{} specifications and the property
being tested.  For the benchmarks in~\autoref{tab:evaluation}, our
results show that writing underapproximate specifications require
anywhere from $6$ - $22$ different operators and include over 100
symbolic events in our most complex examples.  The second group of
columns characterizes \name{}'s ability to discover a violation
compared to the baseline (\textbf{Q2}).  To make the comparison
objective, we execute generators in the two systems 10K times, and
count the number of runs in which a fault was uncovered.  For example,
for \textsf{Stack}, using QCheck to generate random sequences of
$\eff{push}$ and $\eff{pop}$ effects, resulted in an erroneous
execution every 31.9 times on average; in contrast, every execution of
the generator(s) synthesized by \name\ manifested the bug.  Notably,
even for benchmarks in which there is a sophisticated customized
generator available\footnote{We used the artifacts available in their
  publication as the baseline.} (e.g., \textsf{IFC} or \textsf{de
  Brujin}), \tyName{}-guided synthesis leads to significantly better
outcomes.  We use MonkeyDB~\cite{MonkeyDB}, a mock storage system that
can be used to find violations of database isolation properties, as
the baseline for a number of the concurrency benchmarks.  Although it
is specialized for this particular application class, and \name{} is
not, evaluation results are comparable, with \name{} yielding
generators that identify weak isolation violations modestly more
effectively than the baseline.  Similar conclusions apply to the P
benchmarks shown in~\autoref{tab:evaluation2}.  For simple benchmarks,
P's default scheduler, which performs enumerative state exploration to
construct schedules, independently of the behaviors of the actors
under test or the target property, or handwritten ones which inject
additional actors to control input message generation and prevent
uninteresting message orderings (e.g. \textsf{Simplified2PC} or
\textsf{Heartbeat}), are effective and \name{}'s performance is
comparable to these.  As in the previous set of experiments, however,
for more complex benchmarks (e.g., \textsf{ChainReplication},
\textsf{Paxos}), even manual specialization to generate interesting
test sequences is significantly less effective than \name{}'s
trace-guided synthesis approach.  Moreover, as indicated by the
\textsf{Raft} numbers, a random exploration procedure, without
manually engineered guidance, is too na\"{i}ve to find a violating
execution.  It is notable that here too, \name{}'s performance on
these P benchmarks is competitive, and in several cases, superior to
handcrafted P schedulers, even though \name{} has no built-in
knowledge of P's underlying programming model.

The last group of columns in both tables provide details on the cost
of our synthesis procedure. The first column presents total synthesis
time (t$_\text{total}$), which takes anywhere from $.6$ to $165$
seconds (\textbf{Q3}); synthesis time is roughly proportional to
benchmark complexity, as reflected in the \#$\eff{op}$ and \#qualifier
columns.  The last three columns additionally analyze the
characteristics of \name{}'s synthesis procedure with respect to the
number of events in the abstract refined trace (\#evt), the number of
SMT queries (\#SMT), as well as the number of refinement steps
(\#refine); recall that the synthesis procedure uses a refinement loop
to construct a more specialized underapproximation from \tyName{}
specifications that are consistent with the provided safety property.
In general, these numbers correlate with each other, and serve as
rough proxies for benchmark complexity.  They indicate that even for
the most challenging benchmarks, \name{} is able to synthesis an
effective generator in a reasonable amount of time.
\paragraph{Case study}
To demonstrate that \name{} can be effective in real-world scenarios,
we have applied it to \textsf{AnonReadAtomicity}, a distributed
transaction system in use at a major cloud provider, where
asynchronous reads are expected to execute atomically, as discussed
in~\autoref{sec:overview}. The property to be checked requires that if
there exists a key $\Code{k}$ updated within an active transaction,
any successful read response asking its value should return the value
last written to $\Code{k}$ made by that transaction.
%% We can express a
%% violation of this property as the following SRE: %
%% {\small
%%   \begin{align*}
%%     \allA &\msgA{write}{\Code{tid}\;\Code{k}\;\Code{x}\;\Code{OK}}(\anyA\setminus\msgA{write}{\Code{tid}\;\Code{k}\;\Code{x}\;\Code{OK}})^*
%%                 \msgA{readRsp}{\iota\;\Code{tid}\;\Code{k}\;\Code{x}\;\Code{OK}}\;\allA
%% \end{align*}}%
%% \SJ{Fix....}
%% \noindent where the field $\mathit{tid}$ represents a transaction id, while
%% other fields have the same meanings as in the example from
%%\autoref{sec:intro}.
Generating a fault-inducing scenario requires (a) initiating a new
transaction with transaction id $\Code{tid}$, (b) successfully
performing a write within that transaction, and then (c) subsequently
performing a read within $\Code{tid}$ that yields a different value
than the one last written.  As the example in~\ref{sec:overview}
implies, using \tyName{} to express the behavior of $\eff{readRsp}$
enables the synthesis of a test generator that can strategically
request a new transaction to initiate triggering the intended
violation.  A version of the benchmark in which this sequence
structure is enforced manually discovers the violation in 53
executions, while \name{}'s version can uncover the violation on all
of its executions.
